\documentclass{article}
\usepackage{amsmath,ctex,graphicx,booktabs,array,siunitx,wrapfig,caption,subcaption,float,fancyhdr,lastpage,adjustbox,bm,mathtools}
\usepackage[top=1.7cm, bottom=1.5cm, left=1.8cm, right=1.8cm]{geometry}
\usepackage{titling}  
\usepackage{datetime2}  
\pagestyle{fancy}
\fancyhf{}
\fancyhead[C]{\thepage\ / \pageref{LastPage}} 
\renewcommand{\headrulewidth}{0.4pt} 
\pretitle{\begin{center}\LARGE\bfseries}
	\posttitle{\end{center}\vskip 0.2cm}
\newcommand{\crossdot}
{
	\mathbin{\vcenter{\offinterlineskip
			\hbox to 6pt{\hss$\times$\hss}
			\hbox to 6pt{\hss$\cdot$\hss}}}
}
\newcommand{\dotdotdot}
{
	\mathbin{\vcenter{\offinterlineskip
			\hbox to 4pt{\hss$\cdot$\hss}
			\hbox to 4pt{\hss$\cdot$\hss}
			\hbox to 4pt{\hss$\cdot$\hss}}}
}

\begin{document}
	\title{Multipole Expansion of Electromagnetic Field in Cartesian Coordinates}
	\author{Utayah Wu @ School of Physics, Peking University}
	\date{\large Last Updated: \today}
	\maketitle
	\section{Introduction to the Hertz Vector Method}
	
	When a localized electric polarization $\bm{P}$ and magnetization $\bm{M}$ exist in space, and there are no free charges or free currents (which can be neglected in radiation problems since they do not contribute to far-field radiation power), we have
	\[
	\bm{D} = \varepsilon_0 \bm{E} + \bm{P}, \quad 
	\bm{B} = \mu_0 (\bm{H} + \bm{M}), \quad 
	\nabla \cdot \bm{D} = 0, \quad 
	\nabla \times \bm{H} = \frac{\partial \bm{D}}{\partial t}.
	\]
	
	Introducing the scalar potential $\Phi$, the vector potential $\bm{A}$, and adopting the Lorenz gauge, the electromagnetic fields satisfy
	\[
	\bm{E} = -\nabla \Phi - \frac{\partial \bm{A}}{\partial t}, \quad 
	\bm{B} = \nabla \times \bm{A}, \quad 
	\nabla \cdot \bm{A} + \frac{1}{c^2}\frac{\partial \Phi}{\partial t} = 0.
	\]
	
	Substitution yields
	\[
	\left( \nabla^2 - \frac{1}{c^2}\frac{\partial^2}{\partial t^2} \right) \Phi = \frac{1}{\varepsilon_0} \nabla \cdot \bm{P}, \quad 
	\left( \nabla^2 - \frac{1}{c^2}\frac{\partial^2}{\partial t^2} \right) \bm{A} = -\mu_0 \left( \frac{\partial \bm{P}}{\partial t} + \nabla \times \bm{M} \right).
	\]
	
	We now introduce the electric and magnetic Hertz vectors, $\bm{\Pi}_e$ and $\bm{\Pi}_m$, defined by
	\[
	\bm{A} = \mu_0 \left( \frac{\partial \bm{\Pi}_e}{\partial t} + \nabla \times \bm{\Pi}_m \right), \quad 
	\Phi = -\frac{1}{\varepsilon_0} \nabla \cdot \bm{\Pi}_e.
	\]
	
	Substituting these expressions gives
	\[
	\begin{split}
		&\nabla \cdot \left( \nabla^2 \bm{\Pi}_e - \frac{1}{c^2}\frac{\partial^2 \bm{\Pi}_e}{\partial t^2} + \bm{P} \right) = 0, \\
		&\frac{\partial}{\partial t} \left( \nabla^2 \bm{\Pi}_e - \frac{1}{c^2}\frac{\partial^2 \bm{\Pi}_e}{\partial t^2} + \bm{P} \right)
		+ \nabla \times \left( \nabla^2 \bm{\Pi}_m - \frac{1}{c^2}\frac{\partial^2 \bm{\Pi}_m}{\partial t^2} + \bm{M} \right) = 0.
	\end{split}
	\]
	
	A general solution to these equations can be written as (with arbitrary functions $\bm{V}$ and $\xi$):
	\[
	\begin{split}
		\nabla^2 \bm{\Pi}_e - \frac{1}{c^2}\frac{\partial^2 \bm{\Pi}_e}{\partial t^2} &= -\bm{P} + \nabla \times \bm{V}, \\
		\nabla^2 \bm{\Pi}_m - \frac{1}{c^2}\frac{\partial^2 \bm{\Pi}_m}{\partial t^2} &= -\bm{M} - \frac{\partial \bm{V}}{\partial t} + \nabla \left( \frac{\partial \xi}{\partial t} \right).
	\end{split}
	\]
	
	Note that after introducing the Hertz vectors, the scalar and vector potentials still satisfy the Lorenz gauge. The gauge transformation
	\[
	\bm{A} \to \bm{A} + \nabla \Lambda, \quad 
	\Phi \to \Phi - \frac{\partial \Lambda}{\partial t}, \quad 
	\left( \nabla^2 - \frac{1}{c^2}\frac{\partial^2}{\partial t^2} \right) \Lambda = 0,
	\]
	can be expressed in terms of the Hertz vectors as
	\[
	\bm{\Pi}_e \to \bm{\Pi}_e + \nabla \times \bm{G} - \nabla g_e, \quad 
	\bm{\Pi}_m \to \bm{\Pi}_m - \frac{\partial \bm{G}}{\partial t} - \nabla g_m,
	\]
	where $g_m$ and $\bm{G}$ are arbitrary, and $\left( \nabla^2 - \frac{1}{c^2}\frac{\partial^2}{\partial t^2} \right) g_e = 0$.
	Without loss of generality, we may take
	\[
	\left( \nabla^2 - \frac{1}{c^2}\frac{\partial^2}{\partial t^2} \right) \bm{G} = -(\bm{V} + \nabla \xi), \quad g_m = 0.
	\]
	
	Then, the gauge-transformed Hertz vectors satisfy
	\[
	\nabla^2 \bm{\Pi}_e - \frac{1}{c^2}\frac{\partial^2 \bm{\Pi}_e}{\partial t^2} = -\bm{P}, \quad 
	\nabla^2 \bm{\Pi}_m - \frac{1}{c^2}\frac{\partial^2 \bm{\Pi}_m}{\partial t^2} = -\bm{M}.
	\]
	
	Given the spatial distributions of the local polarization and magnetization, the Hertz vectors can be obtained from the above equations. The electromagnetic fields then follow from
	\[
	\begin{split}
		\bm{A} &= \mu_0 \left( \frac{\partial \bm{\Pi}_e}{\partial t} + \nabla \times \bm{\Pi}_m \right), 
		\quad \Phi = -\frac{1}{\varepsilon_0} \nabla \cdot \bm{\Pi}_e, \\[6pt]
		\bm{E} &= -\nabla \Phi - \frac{\partial \bm{A}}{\partial t} 
		= \frac{1}{\varepsilon_0} \nabla (\nabla \cdot \bm{\Pi}_e)
		- \mu_0 \left( \frac{\partial^2 \bm{\Pi}_e}{\partial t^2} + \nabla \times \frac{\partial \bm{\Pi}_m}{\partial t} \right), \\[6pt]
		\bm{B} &= \nabla \times \bm{A} 
		= \mu_0 \left( \nabla \times \frac{\partial \bm{\Pi}_e}{\partial t} + \nabla \times (\nabla \times \bm{\Pi}_m) \right).
	\end{split}
	\]
	
	\section{Derivation of the Primitive Electromagnetic Multipole Moments}
	
	We now derive the explicit relationship between $\bm{P}$, $\bm{M}$, and the local charge distribution.  
	For a field point $\bm{r}$, suppose that the charge is distributed near the point $\bm{r}_s$, which moves with velocity $\bm{v}_s = d\bm{r}_s/dt$. Each charge has a position vector $\bm{r}'$ relative to $\bm{r}_s$, and a corresponding velocity $\bm{v}' = d\bm{r}'/dt$.
	
	Macroscopically, all charges can be regarded as concentrated at $\bm{r}_s$, and the spatial current arises from the motion of $\bm{r}_s$. Microscopically, however, the charge distribution around $\bm{r}_s$ must be considered. Hence,
	\[
	\begin{split}
		\nabla \cdot \bm{E} &= \frac{1}{\varepsilon_0} \int \rho\!\left( \bm{r}' + \bm{r}_s, t \right) \delta\!\left( \bm{r} - \bm{r}' - \bm{r}_s \right) d\bm{r}', \\
		\nabla \cdot \bm{D} &= \int \rho\!\left( \bm{r}' + \bm{r}_s, t \right) \delta\!\left( \bm{r} - \bm{r}_s \right) d\bm{r}', \\
		\nabla \times \bm{B} &= \frac{1}{c^2} \frac{\partial \bm{E}}{\partial t} + \mu_0 \int \rho\!\left( \bm{r}' + \bm{r}_s, t \right) \left( \bm{v}_s + \bm{v}' \right) \delta\!\left( \bm{r} - \bm{r}' - \bm{r}_s \right) d\bm{r}', \\
		\nabla \times \bm{H} &= \frac{\partial \bm{D}}{\partial t} + \int \rho\!\left( \bm{r}' + \bm{r}_s, t \right) \bm{v}_s \delta\!\left( \bm{r} - \bm{r}_s \right) d\bm{r}'.
	\end{split}
	\]
	Here, $d\bm{r}'$ denotes integration over the source volume, and $\delta(\bm{r})$ is the three-dimensional Dirac delta function, whose integral over $d\bm{r}$ equals 1.
	
	Let $\bm{r}^{\otimes n}$ denote the $n$-fold tensor product of $\bm{r}$. A single dot denotes the contraction of the $n$ indices of $\bm{r}^{\otimes n}$ with the first $n$ indices of the subsequent tensor. (Since all tensors that appear later are symmetric in the indices to be contracted, this notation causes no ambiguity.) The operator $\nabla^{\otimes n}$ is defined analogously.
	
	From the Taylor expansion, we have
	\[
	\delta\!\left( \bm{r} - \bm{r}' - \bm{r}_s \right) 
	= \sum_{n=0}^{\infty} \frac{(-1)^n}{n!}\, 
	\bm{r}'^{\otimes n} \cdot \nabla^{\otimes n} \delta\!\left( \bm{r} - \bm{r}_s \right).
	\]
	
	Using $\bm{P} = \bm{D} - \varepsilon_0 \bm{E}$ and $\bm{M} = -\bm{H} + \bm{B}/\mu_0$, we obtain
	\[
	\begin{split}
		\nabla \cdot \bm{P} 
		&= \nabla \cdot \left( \bm{D} - \varepsilon_0 \bm{E} \right) 
		= \sum_{n=1}^{\infty} \frac{(-1)^{n-1}}{n!} 
		\int \rho\!\left( \bm{r}' + \bm{r}_s, t \right) 
		\bm{r}'^{\otimes n} \cdot \nabla^{\otimes n} 
		\delta\!\left( \bm{r} - \bm{r}_s \right) d\bm{r}' \\
		&\to \bm{P} 
		= \sum_{n=1}^{\infty} \frac{(-1)^{n-1}}{n!} 
		\int \rho\!\left( \bm{r}' + \bm{r}_s, t \right) 
		\bm{r}'^{\otimes n} \cdot \nabla^{\otimes (n-1)} 
		\delta\!\left( \bm{r} - \bm{r}_s \right) d\bm{r}', \\
		\nabla \times \bm{M} 
		&= \nabla \times \left( -\bm{H} + \frac{\bm{B}}{\mu_0} \right) 
		= -\frac{\partial \bm{P}}{\partial t} + \bm{v}_s \nabla \cdot \bm{P} 
		+ \sum_{n=1}^{\infty} \frac{(-1)^n}{n!} 
		\int \rho\!\left( \bm{r}' + \bm{r}_s, t \right) 
		\bm{v}' \bm{r}'^{\otimes n} \cdot \nabla^{\otimes n} 
		\delta\!\left( \bm{r} - \bm{r}_s \right) d\bm{r}' \\
		&\to \bm{M} 
		= \sum_{n=1}^{\infty} \frac{(-1)^{n-1}}{n!} \cdot \frac{n}{n+1} 
		\int \rho\!\left( \bm{r}' + \bm{r}_s, t \right) 
		\left( \bm{r}' \times \bm{v}' \right) 
		\bm{r}'^{\otimes (n-1)} \cdot \nabla^{\otimes (n-1)} 
		\delta\!\left( \bm{r} - \bm{r}_s \right) d\bm{r}'.
	\end{split}
	\]
	
	We define the primitive electric and magnetic $2^n$-pole moments as
	\[
	\begin{split}
		P^{(n)}(t) &= \int \rho\!\left( \bm{r}' + \bm{r}_s, t \right) \bm{r}'^{\otimes n} d\bm{r}', \quad
		M^{(n)}(t) = \frac{n}{n+1} \int \bm{r}'^{\otimes (n-1)} \left( \bm{r}' \times \bm{J}\!\left( \bm{r}' + \bm{r}_s, t \right) \right) d\bm{r}',
	\end{split}
	\]
	where $\bm{J}\!\left( \bm{r}' + \bm{r}_s, t \right) = \rho\!\left( \bm{r}' + \bm{r}_s, t \right) \bm{v}'$. Hence, 
	\[
	\begin{split}
		\bm{P}(\bm{r}, t) &= \sum_{n=1}^{\infty} \frac{(-1)^{n-1}}{n!} 
		\nabla^{\otimes (n-1)} \cdot 
		\left( P^{(n)}(t)\, \delta\!\left( \bm{r} - \bm{r}_s \right) \right), \\
		\bm{M}(\bm{r}, t) &= \sum_{n=1}^{\infty} \frac{(-1)^{n-1}}{n!} 
		\nabla^{\otimes (n-1)} \cdot 
		\left( M^{(n)}(t)\, \delta\!\left( \bm{r} - \bm{r}_s \right) \right).
	\end{split}
	\]
	
	According to the Hertz vector formalism,
	\[
	\nabla^2 \bm{\Pi}_e - \frac{1}{c^2}\frac{\partial^2 \bm{\Pi}_e}{\partial t^2} = -\bm{P}, \quad
	\nabla^2 \bm{\Pi}_m - \frac{1}{c^2}\frac{\partial^2 \bm{\Pi}_m}{\partial t^2} = -\bm{M}.
	\]
	
	In an unbounded three-dimensional space, the retarded solutions are
	\[
	\begin{split}
		\bm{\Pi}_e(\bm{r}, t) &= \frac{1}{4\pi} 
		\int \frac{\bm{P}\!\left( \bm{r}', t - \dfrac{|\bm{r} - \bm{r}'|}{c} \right)}{|\bm{r} - \bm{r}'|} d\bm{r}' 
		= \frac{1}{4\pi} \sum_{n=1}^{\infty} \frac{(-1)^{n-1}}{n!} 
		\nabla^{\otimes (n-1)} \cdot \frac{P^{(n)}(t-r/c)}{r}, \\[4pt]
		\bm{\Pi}_m(\bm{r}, t) &= \frac{1}{4\pi} 
		\int \frac{\bm{M}\!\left( \bm{r}', t - \dfrac{|\bm{r} - \bm{r}'|}{c} \right)}{|\bm{r} - \bm{r}'|} d\bm{r}' 
		= \frac{1}{4\pi} \sum_{n=1}^{\infty} \frac{(-1)^{n-1}}{n!} 
		\nabla^{\otimes (n-1)} \cdot \frac{M^{(n)}(t-r/c)}{r}.
	\end{split}
	\]
	
	The corresponding electromagnetic potentials are
	\[
	\begin{split}
		\bm{A} &= \mu_0 \left( \frac{\partial \bm{\Pi}_e}{\partial t} + \nabla \times \bm{\Pi}_m \right) 
		= \frac{\mu_0}{4\pi} \sum_{n=1}^{\infty} \frac{(-1)^{n-1}}{n!} 
		\left( \nabla^{\otimes (n-1)} \cdot \frac{\partial P^{(n)}}{r \partial t} 
		+ \nabla \times \left( \nabla^{\otimes (n-1)} \cdot \frac{M^{(n)}}{r} \right) \right), \\[4pt]
		\Phi &= -\frac{1}{\varepsilon_0} \nabla \cdot \bm{\Pi}_e 
		= -\frac{1}{4\pi\varepsilon_0} \sum_{n=1}^{\infty} \frac{(-1)^{n-1}}{n!} 
		\nabla^{\otimes n} \cdot \frac{P^{(n)}}{r}.
	\end{split}
	\]
	
	This formulation is compact and intuitive. However, note that the n-th radiation field typically involves only the symmetric trace-free (STF) parts of the n-th multipole moments — that is, components that are symmetric under any index exchange and traceless under contraction. (Scalars and vectors are necessarily STF.) Therefore, a mathematical procedure should be applied to convert the above expressions into their STF forms (and compensation terms that can merge with "less than n"-th multipole structures).
	
	
\section{Mathematical Summary of the STF Operation}

We introduce the index-assignment notation $(i_1, \dots, i_n)$.  
For $q_i$ copies of $p_i$-th order symmetric tensors $T^{[i](p_i)}$ ($i = 1, \dots, m$, with $n = \sum_{i=1}^m q_i p_i$),
\[
T^{[1](p_1)}_{(i_1 \dots i_{p_1}} \dots T^{[m](p_m)}_{i_{n-p_m+1} \dots i_n)}
\]
denotes assigning $i_1$ through $i_n$ in all mutually exclusive ways to these $\sum_{i=1}^m q_i$ tensors, and summing over the resulting 
$n!/\!\prod_{i=1}^m \!\big(q_i! (p_i!)^{q_i}\big)$ distinct terms.

For example, $\delta_{(i_1 i_2} f_{i_3)} = \delta_{i_1 i_2} f_{i_3} + \delta_{i_2 i_3} f_{i_1} + \delta_{i_3 i_1} f_{i_2}$.  
As another example, let $T^{(n-2m)}$ be a symmetric tensor of order $n-2m$. Then  
$\delta_{(i_1 i_2} \dots \delta_{i_{2m-3} i_{2m-2}} T^{(n-2m)}_{i_{2m-1} \dots i_{n-2})}$  
corresponds to $q_1 = m-1$, $p_1 = 2$, $q_2 = 1$, $p_2 = n-2m$, and contains  
$n!/\!\big(2^{m-1} (m-1)! (n-2m)!\big)$ distinct terms.

\medskip
Next, we introduce the symmetrization operator $\mathcal{S}$.  
For any $n$-th order tensor $T^{(n)}$, suppose its indices can be divided into $m$ groups, each containing $p_i$ symmetric indices 
($i = 1, \dots, m$, $p_i \ge 1$, and $n - \sum_{i=1}^m p_i$ unsymmetrized). Then
\[
\mathcal{S} T^{(n)} = T_{(i_1 \dots i_n)} \frac{1}{n!} \prod_{i=1}^m p_i!.
\]
When using the index-assignment notation, there are not multiple independent symmetric tensors; rather, one tensor possesses several symmetric index groups (a vector is also a symmetric tensor).  
It is evident that $\mathcal{S}T^{(n)}$ is a fully symmetric tensor of order $n$.

\medskip
The symmetrization operator can also be expressed in the more direct “all permutations” form:
\[
\mathcal{S} T^{(n)} = \frac{1}{n!} \sum_{\text{a.p.}} T_{i_1 \dots i_n},
\]
where “a.p.” (all permutations) denotes all $n!$ possible permutations of $i_1$ through $i_n$. Note that many textbooks' notation "(.)" is equivalent to my $\mathcal{S}$ here, thus multiples of my version of "(.)". 

Since the tensors considered below are typically highly symmetric, explicitly summing over all $n!$ permutations would be cumbersome and unintuitive. Therefore, the shorthand notation $T_{(i_1 \dots i_n)}$ is adopted.  
Note that in some references, $T_{(i_1 \dots i_n)}$ is *defined* to be equivalent to $\mathcal{S}T^{(n)}$ as used here; readers should be aware of this notational distinction.

\medskip
We now introduce the symmetric trace-free (STF) operator $\mathcal{STF}$.  
For any $n$-th order tensor $T^{(n)}$,
\[
\mathcal{STF} T^{(n)} = \mathcal{S} T^{(n)} - \delta_{(i_1 i_2} \Lambda[\mathcal{S} T^{(n)}]_{i_3 \dots i_n)},
\]
where the operator $\Lambda$ lowers the tensor rank by two, and is defined explicitly as
\[
\Lambda[\mathcal{S} T^{(n)}] =
\sum_{m=1}^{\lfloor n/2 \rfloor}
\frac{(-1)^{m-1} (2n - 1 - 2m)!!}{m (2n - 1)!!}
\delta_{(i_1 i_2} \dots \delta_{i_{2m-3} i_{2m-2}}
\mathcal{S} T^{(n;m)}_{i_{2m-1} \dots i_{n-2})}.
\]
Here, $\mathcal{S}T^{(n;m)}$ denotes the tensor obtained by contracting $m$ pairs of indices of the symmetric tensor $\mathcal{S}T^{(n)}$.  
It is straightforward to verify that $\mathcal{STF}T^{(n)}$ is itself an $n$-th order symmetric trace-free tensor.  
Since $\mathcal{STF}$ acts trivially on scalars and vectors (leaving them unchanged), one can also write
\[
\mathcal{S}T^{(n)} = \sum_{m=1}^{\lfloor n/2 \rfloor}
\delta_{(i_1 i_2} \dots \delta_{i_{2m-1} i_{2m}}
\left( \mathcal{STF}\Lambda^{\otimes m}[\mathcal{S}T^{(n)}] \right)_{i_{2m+1} \dots i_n)},
\]
where $\lfloor \cdot \rfloor$ denotes the floor function.
	
	\section{Concrete Steps for STF Operations}
	
	Following the discussion of electric and magnetic polarization, define
	\[
	\begin{split}
		P^{(n)} &= \int \rho\, \bm{r}'^{\otimes n}\, d\bm{r}', \quad
		M^{(n)} = \frac{n}{n+1} \int \bm{r}'^{\otimes (n-1)} \big( \bm{r}' \times \bm{J} \big)\, d\bm{r}' .
	\end{split}
	\]
	Additionally, define
	\[
	\begin{split}
		M^{(n;m)} &= \frac{n}{n+1} \int (r')^{2m}\, \bm{r}'^{\otimes (n-2m-1)} \big( \bm{r}' \times \bm{J} \big)\, d\bm{r}', \\
		N^{(n)} &= \frac{n+1}{n+2} \int \bm{r}'^{\otimes (n-1)} \Big( \bm{r}' \times \big( \bm{r}' \times \bm{J} \big) \Big)\, d\bm{r}', \quad
		N^{(n;m)} = \frac{n+1}{n+2} \int (r')^{2m}\, \bm{r}'^{\otimes (n-2m-1)} \Big( \bm{r}' \times \big( \bm{r}' \times \bm{J} \big) \Big)\, d\bm{r}' .
	\end{split}
	\]
	
	It is straightforward to see that
	\[
	\begin{split}
		\mathcal{S} M^{(n)} &= M^{(n)} - \frac{1}{n} \sum_{p=1}^{n-1} \epsilon_{i_p i_n q}\, N^{(n-1)}_{\,i_1 \dots i_{p-1}\, i_{p+1} \dots i_{n-1}\, q}, \\
		\mathcal{S} N^{(n-1)} &= N^{(n-1)} + \frac{1}{n-1} \sum_{p=1}^{n-2} \epsilon_{i_p i_n q}\, M^{(n;1)}_{\,i_1 \dots i_{p-1}\, i_{p+1} \dots i_{n-2}\, q}.
	\end{split}
	\]
	
	Furthermore (with $\Delta = \nabla^2$):
	\small
	\[
	\begin{split}
		\nabla^{\otimes (n-1)} &\!\cdot\! \left( \frac{P^{(n)}}{r} \right)
		= \nabla^{\otimes (n-1)} \!\cdot\! \left( \frac{\mathcal{STF} P^{(n)}}{r} \right)
		+ (n-1)\, \nabla \!\left( \nabla^{\otimes (n-2)} \!\cdot\! \frac{\Lambda[P^{(n)}]}{r} \right)
		+ \frac{(n-2)(n-1)}{2}\, \Delta\, \nabla^{\otimes (n-3)} \!\cdot\! \left( \frac{\Lambda[P^{(n)}]}{r} \right) \\
		&= \sum_{m=0}^{\left\lfloor \frac{n-1}{2} \right\rfloor} \frac{(n-1)!}{2^m (n-2m-1)!}\,
		\Delta^m\, \nabla^{\otimes (n-2m-1)} \!\cdot\! \frac{\mathcal{STF}\, \Lambda^{\otimes m}[P^{(n)}]}{r}
		\;+\; \nabla \!\left( f_e\, \frac{1}{r} \right)\\
		\nabla^{\otimes (n-1)} &\!\cdot\! \left( \frac{M^{(n)}}{r} \right)
		= \nabla^{\otimes (n-1)} \!\cdot\! \left( \frac{\mathcal{S} M^{(n)}}{r} \right)
		- \frac{n-1}{n}\, \nabla \times \left( \nabla^{\otimes (n-2)} \!\cdot\! \frac{\mathcal{S} N^{(n-1)}}{r} \right) \\
		&\hspace{2cm} + \frac{n-2}{n}\, \Delta\, \nabla^{\otimes (n-3)} \!\cdot\! \left( \frac{\mathcal{S} M^{(n;1)}}{r} \right)
		- \frac{n-2}{n}\, \nabla \!\left( \nabla^{\otimes (n-2)} \!\cdot\! \frac{\mathcal{S} M^{(n;1)}}{r} \right) \\
		&= \sum_{m=0}^{\left\lfloor \frac{n-1}{2} \right\rfloor} \frac{n-2m}{n}\, \Delta^m\, \nabla^{\otimes (n-2m-1)} \!\cdot\! \frac{\mathcal{S} M^{(n;m)}}{r}
		- \nabla \times \sum_{m=0}^{\left\lfloor \frac{n-2}{2} \right\rfloor} \frac{n-2m-1}{n}\, \Delta^m\, \nabla^{\otimes (n-2m-2)} \!\cdot\! \frac{\mathcal{S} N^{(n-1;m)}}{r} + \nabla \!\left( f_m \cdot \frac{1}{r} \right) \\
		&= \sum_{m=0}^{\left\lfloor \frac{n-1}{2} \right\rfloor} \sum_{k=0}^{\left\lfloor \frac{n-2m-1}{2} \right\rfloor}
		\frac{(n-2m)!}{2^k\, n\, (n-2m-1-2k)!}\, \Delta^{m+k}\, \nabla^{\otimes (n-2m-2k-1)} \!\cdot\! \frac{\mathcal{STF}\, \Lambda^{\otimes k}[\mathcal{S} M^{(n;m)}]}{r} \\
		&\quad - \nabla \times \sum_{m=0}^{\left\lfloor \frac{n-2}{2} \right\rfloor} \sum_{k=0}^{\left\lfloor \frac{n-2m-2}{2} \right\rfloor}
		\frac{(n-2m-1)!}{2^k\, n\, (n-2m-2-2k)!}\, \Delta^{m+k}\, \nabla^{\otimes (n-2m-2k-2)} \!\cdot\! \frac{\mathcal{STF}\, \Lambda^{\otimes k}[\mathcal{S} N^{(n-1;m)}]}{r} + \nabla \!\left( f_m\, \frac{1}{r} \right) \\
		&= \sum_{m=0}^{\left\lfloor \frac{n-1}{2} \right\rfloor} \sum_{k=0}^{m}
		\frac{(n-2m+2k)!}{2^k\, n\, (n-2m-1)!}\, \Delta^m\, \nabla^{\otimes (n-2m-1)} \!\cdot\! \frac{\mathcal{STF}\, \Lambda^{\otimes k}[\mathcal{S} M^{(n;m-k)}]}{r} \\
		&\quad - \nabla \times \sum_{m=0}^{\left\lfloor \frac{n-2}{2} \right\rfloor} \sum_{k=0}^{m}
		\frac{(n-2m+2k-1)!}{2^k\, n\, (n-2m-2)!}\, \Delta^m\, \nabla^{\otimes (n-2m-2)} \!\cdot\! \frac{\mathcal{STF}\, \Lambda^{\otimes k}[\mathcal{S} N^{(n-1;m-k)}]}{r}+ \nabla \!\left( f_m\, \frac{1}{r} \right).
	\end{split}
	\]
	\normalsize
	
	Here $f_e$, $\tilde{f}_m$, and $f_m$ are functions of $\nabla$. For functions of this type, $\big( \Delta - \partial^2/(c^2 \partial t^2) \big)\big( f_e/r \big)=0$, so the $f_e$ and $f_m$ terms can be eliminated by a gauge transformation of the Hertz vectors. (In the equalities below, $\stackrel{\text{GT}}{=}$ indicates “gauge transformed.”)
	
	\[
	\begin{split}
		\bm{\Pi}_e &= \frac{1}{4\pi} \sum_{n=1}^{\infty} \frac{(-1)^{n-1}}{n!}\,
		\nabla^{\otimes (n-1)} \!\cdot\! \frac{P^{(n)}}{r} \\
		&\stackrel{\text{GT}}{=} \frac{1}{4\pi} \sum_{n=1}^{\infty} \frac{(-1)^{n-1}}{n!}
		\sum_{m=0}^{\left\lfloor \frac{n-1}{2} \right\rfloor} \frac{(n-1)!}{2^m (n-2m-1)!}\,
		\Delta^m\, \nabla^{\otimes (n-2m-1)} \!\cdot\!
		\frac{\mathcal{STF}\, \Lambda^{\otimes m}[P^{(n)}]}{r}, \\
		\bm{\Pi}_m &= \frac{1}{4\pi} \sum_{n=1}^{\infty} \frac{(-1)^{n-1}}{n!}\,
		\nabla^{\otimes (n-1)} \!\cdot\! \frac{M^{(n)}}{r} \\
		&\stackrel{\text{GT}}{=} \frac{1}{4\pi} \sum_{n=1}^{\infty} \frac{(-1)^{n-1}}{n!}
		\sum_{m=0}^{\left\lfloor \frac{n-1}{2} \right\rfloor} \sum_{k=0}^{m}
		\frac{(n-2m+2k)!}{2^k (n-2m-1)!}\, \Delta^m \Bigg[
		\frac{1}{n}\, \nabla^{\otimes (n-2m-1)} \!\cdot\!
		\frac{\mathcal{STF}\, \Lambda^{\otimes k}[\mathcal{S} M^{(n;m-k)}]}{r} \\
		&\qquad + \frac{1}{(n+1)^2}\,
		\nabla \times \left( \nabla^{\otimes (n-2m-1)} \!\cdot\!
		\frac{\mathcal{STF}\, \Lambda^{\otimes k}[\mathcal{S} N^{(n;m-k)}]}{r} \right) \Bigg].
	\end{split}
	\]
	
	In the derivations above we used the rearrangements
	\[
	\sum_{m=0}^{\left\lfloor \frac{n-2}{2} \right\rfloor} \sum_{k=0}^{\left\lfloor \frac{n-2m-2}{2} \right\rfloor} f(m,k)
	= \sum_{m=0}^{\left\lfloor \frac{n-2}{2} \right\rfloor} \sum_{k=0}^{m} f(m-k,k), \quad
	\sum_{n=1}^{\infty} \sum_{m=0}^{\left\lfloor \frac{n-1}{2} \right\rfloor} x^{\,n-2m-1}\, f(n,m)
	= \sum_{n=1}^{\infty} \sum_{m=0}^{\infty} x^{\,n-1}\, f(n+2m,m).
	\]
	
	Perform the gauge transformation $\bm{\Pi}_e \to \bm{\Pi}_e + \nabla \times \bm{G} - \nabla g_e$, $\bm{\Pi}_m \to \bm{\Pi}_m - \partial \bm{G}/\partial t - \nabla g_m$. Introduce the gauge functions with $g_m = 0$,
	\[
	\begin{split}
		\bm{G} &= \frac{1}{4\pi} \cdot \frac{\partial}{c^2 \partial t} \nabla \times \sum_{n=1}^{\infty} \frac{(-1)^{n-1}}{n!} \sum_{m=0}^{\left\lfloor \frac{n-1}{2} \right\rfloor} \sum_{k=0}^m \frac{(n-2m+2k)!}{2^k (n-2m-1)!} \Delta^{m-1} \cdot \frac{1}{(n+1)^2} \nabla \times \left( \nabla^{\otimes (n-2m-1)} \cdot \frac{\mathcal{STF}\Lambda^{\otimes k}[\mathcal{S}N^{(n;m-k)}]}{r} \right), \\
		g_e &= \frac{1}{4\pi} \cdot \frac{\partial}{c^2 \partial t} \sum_{n=1}^{\infty} \frac{(-1)^{n-1}}{n!} \sum_{m=0}^{\left\lfloor \frac{n-1}{2} \right\rfloor} \sum_{k=0}^m \frac{(n-2m+2k)!}{2^k (n-2m-1)!} \Delta^{m-1} \cdot \frac{1}{(n+1)^2} \nabla^{\otimes (n-2m)} \cdot \frac{\mathcal{STF}\Lambda^{\otimes k}[\mathcal{S}N^{(n;m-k)}]}{r}.
	\end{split}
	\]
	
	It is readily seen that $\left( \Delta - \partial^2/(c^2 \partial t^2) \right) g_e = 0$, and thus
	\[
	\bm{\Pi}_e \stackrel{\text{GT}}{=} \frac{1}{4\pi} \sum_{n=1}^{\infty} \frac{(-1)^{n-1}}{n!} \nabla^{\otimes (n-1)} \cdot \frac{\bm{\mathcal{P}}^{(n)}}{r}, \quad
	\bm{\Pi}_m \stackrel{\text{GT}}{=} \frac{1}{4\pi} \sum_{n=1}^{\infty} \frac{(-1)^{n-1}}{n!} \nabla^{\otimes (n-1)} \cdot \frac{\bm{\mathcal{M}}^{(n)}}{r}.
	\]
	
	The STF-processed electromagnetic multipole moments are
	\[
	\begin{split}
		\bm{\mathcal{P}}^{(n)} &= \sum_{m=0}^{\infty} \Delta^m \left( \frac{n}{2^m (n+2m)} \mathcal{STF}\Lambda^{\otimes m}[P^{(n+2m)}]
		- \sum_{k=0}^m \frac{n (n+2k)!}{2^k (n+2m+1) (n+2m+1)!} \mathcal{STF}\Lambda^{\otimes k}\!\left[ \mathcal{S} \frac{\partial N^{(n+2m;m-k)}}{c^2 \partial t} \right] \right), \\
		\bm{\mathcal{M}}^{(n)} &= \sum_{m=0}^{\infty} \Delta^m \sum_{k=0}^m \frac{n (n+2k)!}{2^k (n+2m) (n+2m)!} \mathcal{STF}\Lambda^{\otimes k}[\mathcal{S}M^{(n+2m;m-k)}].
	\end{split}
	\]
	
	The first few terms (here one may view $\Delta$ as $\partial^2/(c^2 \partial t^2)$ at the level of functional dependence) are
	\[
	\begin{split}
		&\bm{\mathcal{P}}^{(n)} = \mathcal{STF}P^{(n)}
		+ \frac{\partial}{c^2 \partial t} \left( \frac{n}{2(n+2)} \mathcal{STF}\Lambda\!\left[ \frac{\partial P^{(n+2)}}{\partial t} \right]
		- \frac{n}{(n+1)^2} \mathcal{STF}N^{(n)} \right) \\
		&\quad + \frac{\partial^3}{c^4 \partial t^3} \left( \frac{n}{4(n+4)} \mathcal{STF}\Lambda^{\otimes 2}\!\left[ \frac{\partial P^{(n+4)}}{\partial t} \right]
		- \frac{n}{2(n+1)(n+2)(n+3)^2} \mathcal{STF}N^{(n+2;1)}
		- \frac{n}{2(n+3)^2} \mathcal{STF}\Lambda[\mathcal{S}N^{(n+2)}] \right) + \dots, \\
		&\bm{\mathcal{M}}^{(n)} = \mathcal{STF}M^{(n)}
		+ \frac{\partial}{c^2 \partial t^2} \left( \frac{n}{2(n+2)} \mathcal{STF}\Lambda[\mathcal{S}M^{(n+2)}]
		+ \frac{n}{(n+2)^2(n+1)} \mathcal{STF}M^{(n+2;1)} \right) + \dots.
	\end{split}
	\]
	
	Mathematically, $\mathcal{S}\mathcal{T}\mathcal{F}\, P^{(n)}$ and $\mathcal{S}\mathcal{T}\mathcal{F}\, M^{(n)}$ can be written in more compact forms:
	\[
	\begin{split}
		\mathcal{S}\mathcal{T}\mathcal{F}\, P^{(n)} &= \mathcal{S}\mathcal{T}\mathcal{F} \int \rho \, \bm{r}'^{\otimes n} \, d\bm{r}'
		= \frac{(-1)^n}{(2n-1)!!} \int \rho \, (r')^{2n+1} \, \nabla'^{\otimes n} \frac{1}{r'} \, d\bm{r}', \\
		\mathcal{S}\mathcal{T}\mathcal{F}\, M^{(n)} &=\mathcal{S}\mathcal{T}\mathcal{F}\, \frac{n}{n+1} \int \bm{r}'^{\otimes (n-1)} \left( \bm{r}' \times \bm{J} \right) d\bm{r}'
		= \frac{(-1)^{n+1}}{(n+1)(2n-1)!!} \int (r')^{2n+1} \, (\bm{J} \times \nabla')_{(i_1} \nabla'_{i_2} \cdots \nabla'_{i_n)} \frac{1}{r'} \, d\bm{r}'.
	\end{split}
	\]
	
	From charge conservation, $\partial \rho/\partial t + \nabla' \cdot \bm{J} = 0$, and by Gauss's theorem, we have
	\[
	\frac{\partial P^{(n)}}{\partial t} = \int r'_{i_1} \dots r'_{i_{n-1}} J_{i_n} \, d\bm{r}'.
	\]
	
	Thus $\partial P^{(n)}/\partial t$, $M^{(n-1)}$, and $N^{(n-2)}$ all share the same dimensions as $\int (r')^{n-1} J\, d\bm{r}'$, and belong to the same order in the expansion. Standard electrodynamics textbooks usually expand only up to the order $\int r' J\, d\bm{r}'$; most papers involving toroidal multipole moments expand to the order $\int (r')^3 J\, d\bm{r}'$; expansions to the order $\int (r')^4 J\, d\bm{r}'$ are extremely rare.
	
	Following the method above, the expansion terms at that order should be
	\[
	\begin{split}
		\bm{\mathcal{P}}^{(1)} &= \int \bm{r}' \rho \, d\bm{r}' \;-\; \frac{\partial}{c^2 \partial t} \cdot \frac{1}{10} \int \!\left( (\bm{r}' \cdot \bm{J}) \bm{r}' - 2 (r')^2 \bm{J} \right) d\bm{r}'
		+ \frac{\partial^3}{c^4 \partial t^3} \cdot \frac{1}{1680} \int \!\left( 11 (r')^4 \bm{J} - 5 (\bm{r}' \cdot \bm{J}) (r')^2 \bm{r}' \right) d\bm{r}', \\
		\bm{\mathcal{P}}^{(2)} &= \int \!\left( \bm{r}' \bm{r}' - \frac{1}{3} (r')^2 \bm{I} \right) \rho \, d\bm{r}'
		- \frac{\partial}{c^2 \partial t} \cdot \frac{1}{42} \int \!\left( 4 \bm{r}' \bm{r}' (\bm{r}' \cdot \bm{J}) - 5 (r')^2 (\bm{r}' \bm{J} + \bm{J} \bm{r}')
		+ 2 (r')^2 (\bm{r}' \cdot \bm{J}) \bm{I} \right) d\bm{r}', \\
		\bm{\mathcal{P}}^{(3)} &= \int \!\left( \bm{r}' \bm{r}' \bm{r}' - \frac{1}{5} (r')^2 \delta_{(i_1 i_2} r'_{i_3)} \right) \rho \, d\bm{r}'
		- \frac{\partial}{c^2 \partial t} \cdot \frac{1}{60} \int \Bigg[ 5 \bm{r}' \bm{r}' \bm{r}' (\bm{r}' \cdot \bm{J})
		- 5 (r')^2 r'_{(i_1} r'_{i_2} J_{i_3)} \\
		&\qquad+ (r')^2 (\bm{r}' \cdot \bm{J}) \delta_{(i_1 i_2} r'_{i_3)}
		+ (r')^4 \delta_{(i_1 i_2} J_{i_3)} \Bigg] d\bm{r}', \\
		\bm{\mathcal{P}}^{(4)} &= \int \!\left( \bm{r}' \bm{r}' \bm{r}' \bm{r}' - \frac{1}{7} (r')^2 \delta_{(i_1 i_2} r'_{i_3} r'_{i_4)}
		+ \frac{1}{35} (r')^4 \delta_{(i_1 i_2} \delta_{i_3 i_4)} \right) \rho \, d\bm{r}', \\
		\bm{\mathcal{P}}^{(5)} &= \int \!\left( \bm{r}'^{\otimes 5} - \frac{1}{9} (r')^2 \delta_{(i_1 i_2} r'_{i_3} r'_{i_4} r'_{i_5)}
		+ \frac{1}{63} (r')^4 \delta_{(i_1 i_2} \delta_{i_3 i_4} r'_{i_5)} \right) \rho \, d\bm{r}', \\
		\bm{\mathcal{M}}^{(1)} &= \frac{1}{2} \int \bm{r}' \times \bm{J} \, d\bm{r}' + \frac{1}{20} \int (r')^2 (\bm{r}' \times \bm{J}) \, d\bm{r}', \\
		\bm{\mathcal{M}}^{(2)} &= \frac{1}{3} \int \left( \bm{r}' (\bm{r}' \times \bm{J}) + (\bm{r}' \times \bm{J}) \bm{r}' \right) d\bm{r}'
		+ \frac{1}{42} \int (r')^2 \left( \bm{r}' (\bm{r}' \times \bm{J}) + (\bm{r}' \times \bm{J}) \bm{r}' \right) d\bm{r}', \\
		\bm{\mathcal{M}}^{(3)} &= \frac{1}{4} \int \left( \bm{r}' \bm{r}' (\bm{r}' \times \bm{J}) + \bm{r}' (\bm{r}' \times \bm{J}) \bm{r}' + (\bm{r}' \times \bm{J}) \bm{r}' \bm{r}' - \frac{1}{5} (r')^2 \delta_{(i_1 i_2} (\bm{r}' \times \bm{J})_{i_3)} \right) d\bm{r}', \\
		\bm{\mathcal{M}}^{(4)} &= \frac{1}{5} \int \left( r'_{(i_1} r'_{i_2} r'_{i_3} (\bm{r}' \times \bm{J})_{i_4)} - \frac{1}{7} \delta_{(i_1 i_2} \left( r'_{i_3} (\bm{r}' \times \bm{J})_{i_4)} + (\bm{r}' \times \bm{J})_{i_3} r'_{i_4)} \right) \right) d\bm{r}'.
	\end{split}
	\]
	
	\section{Electromagnetic Fields of STF Multipole Moments}
	
	Once the electromagnetic multipole moments have been processed into STF form, the radiation fields can be written compactly. From the Hertz vectors, the fields follow as
	\[
	\begin{split}
		\bm{E} &= \frac{1}{\varepsilon_0} \nabla (\nabla \cdot \bm{\Pi}_e)
		- \mu_0 \left( \frac{\partial^2 \bm{\Pi}_e}{\partial t^2} + \nabla \times \frac{\partial \bm{\Pi}_m}{\partial t} \right) \\
		&= \frac{1}{4\pi \varepsilon_0} \sum_{n=1}^{\infty} \frac{(-1)^{n-1}}{n!}
		\nabla \times \left( \nabla \times \left( \nabla^{\otimes (n-1)} \cdot \frac{\bm{\mathcal{P}}^{(n)}}{r} \right) \right)
		- \frac{\mu_0}{4\pi} \sum_{n=1}^{\infty} \frac{(-1)^{n-1}}{n!}
		\nabla \times \left( \nabla^{\otimes (n-1)} \cdot \frac{\partial \bm{\mathcal{M}}^{(n)}}{r \partial t} \right), \\
		\bm{B} &= \mu_0 \left( \nabla \times \frac{\partial \bm{\Pi}_e}{\partial t}
		+ \nabla \times (\nabla \times \bm{\Pi}_m) \right) \\
		&= \frac{\mu_0}{4\pi} \sum_{n=1}^{\infty} \frac{(-1)^{n-1}}{n!}
		\nabla \times \left( \nabla^{\otimes (n-1)} \cdot \frac{\partial \bm{\mathcal{P}}^{(n)}}{r \partial t} \right)
		+ \frac{\mu_0}{4\pi} \sum_{n=1}^{\infty} \frac{(-1)^{n-1}}{n!}
		\nabla \times \left( \nabla \times \left( \nabla^{\otimes (n-1)} \cdot \frac{\bm{\mathcal{M}}^{(n)}}{r} \right) \right).
	\end{split}
	\]
	
	Keeping only the contributions relevant to energy and angular-momentum radiation, we obtain
	\[
	\begin{split}
		\bm{E}_{\text{rad}} &=
		\frac{1}{4\pi \varepsilon_0 r} \sum_{n=1}^{\infty} \frac{1}{n!}
		\left[ \left( \left( \bm{\hat{r}}^{\otimes (n-1)} \cdot \frac{\partial^{n+1} \bm{\mathcal{P}}^{(n)}}{c^{n+1} \partial t^{n+1}} \right) \times \bm{\hat{r}} \right) \times \bm{\hat{r}}
		- \frac{1}{c} \left( \bm{\hat{r}}^{\otimes (n-1)} \cdot \frac{\partial^{n+1} \bm{\mathcal{M}}^{(n)}}{c^{n+1} \partial t^{n+1}} \right) \times \bm{\hat{r}} \right] \\
		&\quad + \frac{1}{4\pi \varepsilon_0 r^2} \sum_{n=1}^{\infty} \frac{1}{n!} \cdot \frac{n+1}{2}
		\left( (n+2)\, \bm{\hat{r}} \left( \bm{\hat{r}}^{\otimes n} \cdot \frac{\partial^n \bm{\mathcal{P}}^{(n)}}{c^n \partial t^n} \right)
		- n \left( \bm{\hat{r}}^{\otimes (n-1)} \cdot \frac{\partial^n \bm{\mathcal{P}}^{(n)}}{c^n \partial t^n} \right) \right), \\
		\bm{B}_{\text{rad}} &=
		\frac{\mu_0}{4\pi r} \sum_{n=1}^{\infty} \frac{1}{n!}
		\left[ \left( \left( \bm{\hat{r}}^{\otimes (n-1)} \cdot \frac{\partial^{n+1} \bm{\mathcal{M}}^{(n)}}{c^{n+1} \partial t^{n+1}} \right) \times \bm{\hat{r}} \right) \times \bm{\hat{r}}
		+ c \left( \bm{\hat{r}}^{\otimes (n-1)} \cdot \frac{\partial^{n+1} \bm{\mathcal{P}}^{(n)}}{c^{n+1} \partial t^{n+1}} \right) \times \bm{\hat{r}} \right] \\
		&\quad + \frac{\mu_0}{4\pi r^2} \sum_{n=1}^{\infty} \frac{1}{n!} \cdot \frac{n+1}{2}
		\left( (n+2)\, \bm{\hat{r}} \left( \bm{\hat{r}}^{\otimes n} \cdot \frac{\partial^n \bm{\mathcal{M}}^{(n)}}{c^n \partial t^n} \right)
		- n \left( \bm{\hat{r}}^{\otimes (n-1)} \cdot \frac{\partial^n \bm{\mathcal{M}}^{(n)}}{c^n \partial t^n} \right) \right).
	\end{split}
	\]
	
	The instantaneously radiated power is
	\small
	\[
	\begin{split}
		\frac{d^2 E_{\text{rad}}}{dt\,d\Omega}
		&= \left. \frac{r^2}{\mu_0} \, (\bm{E} \times \bm{B}) \cdot \bm{\hat{r}} \right|_{r \to \infty}
		= r^2 \varepsilon_0 c \, |\bm{E}|^2_{r \to \infty} \\
		&= \frac{c}{16\pi^2 \varepsilon_0} \sum_{n=1}^{\infty} \frac{1}{(n!)^2} \left(
		\left| \left( \bm{\hat{r}}^{\otimes (n-1)} \cdot \frac{\partial^{n+1} \bm{\mathcal{P}}^{(n)}}{c^{n+1} \partial t^{n+1}} \right) \times \bm{\hat{r}} \right|^2
		+ \left| \left( \bm{\hat{r}}^{\otimes (n-1)} \cdot \frac{\partial^{n+1} \bm{\mathcal{M}}^{(n)}}{c^{n+2} \partial t^{n+1}} \right) \times \bm{\hat{r}} \right|^2 \right)
		+ \text{ct}, \\
		\frac{dE_{\text{rad}}}{dt} &= \frac{c}{4\pi \varepsilon_0} \sum_{n=1}^{\infty} \frac{2^n (n+1)}{n (2n+1)!}
		\left( \left\| \frac{\partial^{n+1} \bm{\mathcal{P}}^{(n)}}{c^{n+1} \partial t^{n+1}} \right\|^2
		+ \left\| \frac{\partial^{n+1} \bm{\mathcal{M}}^{(n)}}{c^{n+2} \partial t^{n+1}} \right\|^2 \right).
	\end{split}
	\]
	\normalsize
	
	The instantaneously radiated angular-momentum is
	\small
\[
\begin{split}
	\frac{d^2 \bm{L}_{\text{rad}}}{dt \, d\Omega}
	&= \varepsilon_0 \left. r^3 \left( (\bm{E} \cdot \bm{\hat{r}})(\bm{E} \times \bm{\hat{r}})
	+ c^2 (\bm{B} \cdot \bm{\hat{r}})(\bm{B} \times \bm{\hat{r}}) \right) \right|_{r \to \infty} \\
	&= \frac{1}{16\pi^2 \varepsilon_0} \sum_{n=1}^{\infty} \frac{n+1}{(n!)^2} \left[
	\bm{\hat{r}} \times \left( \bm{\hat{r}}^{\otimes (n-1)} \cdot \frac{\partial^{n+1} \bm{\mathcal{P}}^{(n)}}{c^{n+1} \partial t^{n+1}} \right)
	\left( \bm{\hat{r}}^{\otimes n} \cdot \frac{\partial^n \bm{\mathcal{P}}^{(n)}}{c^n \partial t^n} \right) + \bm{\hat{r}} \times \left( \bm{\hat{r}}^{\otimes (n-1)} \cdot \frac{\partial^{n+1} \bm{\mathcal{M}}^{(n)}}{c^{n+2} \partial t^{n+1}} \right)
	\left( \bm{\hat{r}}^{\otimes n} \cdot \frac{\partial^n \bm{\mathcal{M}}^{(n)}}{c^{n+1} \partial t^n} \right) \right] + \text{ct}, \\
	\frac{d \bm{L}_{\text{rad}}}{dt}
	&= \frac{1}{4\pi \varepsilon_0} \sum_{n=1}^{\infty} \frac{2^n (n+1)}{(2n+1)!} \left(
	\frac{\partial^{n+1} \bm{\mathcal{P}}^{(n)}}{c^{n+1} \partial t^{n+1}} \crossdot
	\frac{\partial^n \bm{\mathcal{P}}^{(n)}}{c^n \partial t^n}
	+ \frac{\partial^{n+1} \bm{\mathcal{M}}^{(n)}}{c^{n+2} \partial t^{n+1}} \crossdot
	\frac{\partial^n \bm{\mathcal{M}}^{(n)}}{c^{n+1} \partial t^n} \right).
\end{split}
\]
	\normalsize
	
	Here, the cross-terms (ct) arising from couplings between different multipole orders vanish upon integration over solid angle. The notation $\|\cdot\|^2$ denotes the sum of squares over all tensor components, and $\crossdot$ denotes taking the cross product over one index while contracting over the remaining indices. In the calculations we use
	\[
	\int \hat{r}_{i_1} \cdots \hat{r}_{i_{2n}}\, d\Omega
	= \frac{4\pi}{(2n+1)!!}\, \delta_{(i_1 i_2} \cdots \delta_{i_{2n-1} i_{2n})}.
	\]
	
	It is interesting to note that if we formally express $\partial/\partial t$ as $-i\omega$ and neglect the distinction between $\cdot$ and $\crossdot$, the $n$-th radiative multipole moment satisfies $L_{\text{rad}}^{(n)}/E_{\text{rad}}^{(n)} = n/\omega$. This result appears to imply a connection between the $n$-th multipole and a particle with spin quantum number $n$—the latter also has an angular momentum-to-energy ratio of $L_z/E = n\hbar/(\hbar\omega) = n/\omega$.
	
\section{Comparison with Known Results}

In the following paper, the authors expanded their analysis up to the order of $\int (r')^4 J \, d\bm{r}'$. Up to $\int (r')^3 J \, d\bm{r}'$, their results agree perfectly with mine; however, at the $\int (r')^4 J \, d\bm{r}'$ correction order for the electric dipole and electric octupole moments, a significant discrepancy arises.  

The paper is:  

E. A. Gurvitz, K. S. Ladutenko, P. A. Dergachev, A. B. Evlyukhin, A. E. Miroshnichenko, and A. S. Shalin,  
$\textit{“The high-order toroidal moments and anapole states in all-dielectric photonics,”}$arXiv:1810.04945, Oct.2018.

Specifically, let unadorned symbols denote my results, a tilde ($\tilde{\phantom{a}}$) denote those from the paper, and a superscript on the left indicate the correction order:
\[
\begin{split}
	\bm{\prescript{2}{}{\mathcal{P}}}^{(1)} &= \frac{1}{1680} \int \!\left( 11 (r')^4 \bm{J} - 5 (\bm{r}' \cdot \bm{J}) (r')^2 \bm{r}' \right) d\bm{r}', \\
	\bm{\prescript{2}{}{\tilde{\mathcal{P}}}}^{(1)} &= \frac{1}{280} \int \!\left( 3 (r')^4 \bm{J} - 2 (\bm{r}' \cdot \bm{J}) (r')^2 \bm{r}' \right) d\bm{r}', \\
	\bm{\prescript{1}{}{\mathcal{P}}}^{(3)} &= \frac{1}{60} \int \Bigg[ 5 \bm{r}' \bm{r}' \bm{r}' (\bm{r}' \cdot \bm{J})
	- 5 (r')^2 r'_{(i_1} r'_{i_2} J_{i_3)} + (r')^2 (\bm{r}' \cdot \bm{J}) \delta_{(i_1 i_2} r'_{i_3)}
	+ (r')^4 \delta_{(i_1 i_2} J_{i_3)} \Bigg] d\bm{r}', \\
	\bm{\prescript{1}{}{\tilde{\mathcal{P}}}}^{(3)} &= \frac{1}{300} \int \Bigg[ 35 \bm{r}' \bm{r}' \bm{r}' (\bm{r}' \cdot \bm{J})
	- 20 (r')^2 r'_{(i_1} r'_{i_2} J_{i_3)} + (r')^2 (\bm{r}' \cdot \bm{J}) \delta_{(i_1 i_2} r'_{i_3)}
	+ 4 (r')^4 \delta_{(i_1 i_2} J_{i_3)} \Bigg] d\bm{r}'.
\end{split}
\]

Their differences are
\[
\begin{split}
	\bm{\prescript{2}{}{\mathcal{P}}}^{(1)} - \bm{\prescript{2}{}{\tilde{\mathcal{P}}}}^{(1)} &=
	\frac{1}{240} \int \!\left( - (r')^4 \bm{J} + (\bm{r}' \cdot \bm{J}) (r')^2 \bm{r}' \right) d\bm{r}', \\
	\bm{\prescript{1}{}{\mathcal{P}}}^{(3)} - \bm{\prescript{1}{}{\tilde{\mathcal{P}}}}^{(3)} &=
	\frac{1}{300} \int \Bigg[ -10 \bm{r}' \bm{r}' \bm{r}' (\bm{r}' \cdot \bm{J})
	- 5 (r')^2 r'_{(i_1} r'_{i_2} J_{i_3)} + 4 (r')^2 (\bm{r}' \cdot \bm{J}) \delta_{(i_1 i_2} r'_{i_3)}
	+ (r')^4 \delta_{(i_1 i_2} J_{i_3)} \Bigg] d\bm{r}'.
\end{split}
\]

Intuitively, since the formulation based on fully STF-processed electromagnetic multipole moments should be unique, this discrepancy is rather puzzling.

\end{document}